\documentclass[11pt]{article}
\usepackage{amsmath,amssymb,amsthm,hyperref}
\newtheorem{theorem}{Theorem}
\newtheorem{lemma}[theorem]{Lemma}
\newcommand{\Q}{\mathbb Q}
\newcommand{\m}{\mathfrak m}
\title{Complete finite reduction of mixed-degree GST curvature folds}
\author{Hand mathematics before Lean; author self-review}
\date{3 November 2026}
\begin{document}\maketitle

\section{Target and precise category}
Start at \texttt{fde04bf40c9482f6818521a04df3700b2edb4210} with arbitrary
smooth projective complex $X$, positive codimension $p$, and prescribed
$\alpha\in\operatorname{Hdg}^p(X)_\Q$. The universal target is a rational
Chow cycle on this original $X$ with canonical Betti class $\alpha$.
Work componentwise on connected $X$. Put
\[
 A=\bigoplus_{q=0}^{n}\operatorname{Hdg}^q(X)_\Q,\qquad
 \m=A_{>0},\qquad \m^{n+1}=0,\qquad A/\m=\Q.
\]
The weight grading $q$ means actual cohomological degree $2q$.
Let $\alpha_i\in A_{p_i}$, $p_i\geq1$, be any finite list, including
the prescribed input. On the actual product of rational homotopy-fiber
torsors over $X$, retain the previously constructed model
\[
 B=H^*(X,\R)\otimes\Lambda(u_i,v_i),\quad
 |u_i|=|v_i|=2p_i-1,\quad du_i=\alpha_i,\quad dv_i=0.                 \tag{1}
\]
Its conjugate generators $\beta_i=u_i+iv_i$ have specified type
$(p_i,p_i-1)$. This is the model filtration, not a claimed compact complex
realization of a higher rational torsor.

Allowed changes are homogeneous invertible matrices over $A$, applied
to both the $u$ and $v$ lists. Entry $M_{ji}$ has weight $p_j-p_i$;
negative weights are zero. They preserve (1), conjugation, and the model
filtration. This is a precise operation category. The theorem below does
not classify arbitrary nonlinear changes involving products of odd handles.

\section{A terminating algorithm and the sharp minimum}
Set $I=(\alpha_1,\ldots,\alpha_r)\subset A$ and
\[
 \mu_q=\dim_\Q(I/\m I)_q,\qquad \mu=\sum_q\mu_q.                  \tag{2}
\]
These are ideal generators, not the dimension of $I$ and not the rank of
the original curvature list as a rational vector space.

\begin{theorem}[Complete mixed-degree curvature normal form]
Every list in (1) can be transformed by finitely many allowed changes to
exactly $\mu_q$ nonzero curvature handles of weight $q$, and
$r_q-\mu_q$ zero-curvature handle pairs of that weight, where
$r_q=\#\{i:p_i=q\}$. No allowed change can leave fewer nonzero curvature
handles, in any weight. Two minimal generating lists for the same ideal
are related by a homogeneous invertible matrix. The active curvatures are
not claimed to be a canonical ordered basis.
\end{theorem}
\begin{proof}
Process weights in ascending order. At weight $q$, form the subspace
$L_q=\sum_{s<q}A_{q-s}I_s$. It is computed from already processed lower
weights and equals $(\m I)_q$. The images of the original weight-$q$
curvatures span $I_q/L_q$: all other contributions to $I_q$ have smaller
curvature weight. Choose from those images a basis by rational Gaussian
elimination. Its size is $\mu_q$. An invertible rational change within
weight $q$ makes every remaining curvature lie in $L_q$. Express each
such remainder as a homogeneous $A$-linear combination of the retained
lower-weight curvatures. This is possible inductively because lower-weight
curvatures already generate $I$ below $q$. Subtract that combination from
the corresponding $u$ and $v$ generators. This exposes a closed pair.
The diagonal changes are invertible over $\Q$ and all other changes are
strictly triangular in weight. Therefore the entire matrix is invertible.
The procedure ends after the finitely many listed weights.

For any allowed presentation, the surviving curvature generators must
span $I/\m I$. A zero-curvature row contributes zero there. Its weight-$q$
part thus requires at least $\mu_q$ active rows. This proves optimality.

If $a$ and $b$ are minimal lists generating $I$, write $b=Ma$ with
homogeneous entries. Their images are bases of $I/\m I$, so every
equal-weight diagonal block of $M$ is an invertible rational matrix.
The matrix is block triangular in weight; its inverse is the finite
triangular substitution. Thus it is an allowed change. No choice of
division by a positive-degree class has been made.
\end{proof}

Equivalently, the standard nilpotent-ideal argument is constructive here.
If a submodule $J\subset I$ maps onto $I/\m I$, then
$I=J+\m I=J+\m^2I=\cdots=J+\m^{n+1}I=J$.
This supplies a finite bound on propagation; it does not require a
Noetherian completion, a regular sequence, or a nonzero determinant of
an intersection matrix.

\section{Full cohomology factor and actual higher torsors}
\begin{theorem}[Complete split of every redundant pair]
There is a conjugation-compatible filtered DGA isomorphism over the base
\[
 B\cong B_{\min}\otimes
 \Lambda\bigl(w_{q,j},z_{q,j}:1\leq j\leq r_q-\mu_q\bigr),
 \quad dw_{q,j}=dz_{q,j}=0,\quad |w_{q,j}|=|z_{q,j}|=2q-1.            \tag{3}
\]
Consequently the full Betti polynomial satisfies
\[
 P_B(t)=P_{B_{\min}}(t)\prod_q(1+t^{2q-1})^{2(r_q-\mu_q)}.           \tag{4}
\]
The same substitution is an equivalence over $X$ of the actual rational
topological torsor constructions, after choices of classifying maps and
homotopies. It is not a new principal elliptic quotient in higher degree.
\end{theorem}
\begin{proof}
Every row subtraction sends $u_j$ to $u_j-\sum_i a_i u_i$ and $v_j$ to
$v_j-\sum_i a_i v_i$. Its inverse adds the same expression; the unchanged
rows make the row operator square zero. Even closed coefficients ensure
the chain equation and the absence of an extra sign. Hodge types agree
because $a_i$ has type $(p_j-p_i,p_j-p_i)$. For a zero curvature row,
both generators have zero differential and their conjugate combinations
are pure. Splitting exterior generators proves (3); the field Kunneth
formula proves (4) in every degree, including an empty list.

For actual torsors, use the path-space construction from the preceding
derivation. Multiplication by a base class is implemented by a cup-product
map of Eilenberg--Mac Lane spaces. Subtract the corresponding cup paths
from the target nullhomotopy, retaining every source path. Adding them
is the inverse up to the chosen homotopies. Equal-weight rational changes
are invertible maps of rational Eilenberg--Mac Lane groups. A zero
classifying class has a chosen nullhomotopy and hence a product fiber.
These constructions require no nilpotence hypothesis on $\pi_1(X)$.
\end{proof}
The coefficient algebra in (2) is the rational Hodge subalgebra, whereas
$B_{\min}$ retains \emph{all} of $H^*(X,\R)$, including odd and off-diagonal
classes. Formula (4) is not obtained by discarding the rest of the base.

\section{The prescribed input cannot be erased by its own cup births}
\begin{theorem}[Exact single-target conservation under cup amplification]
Take $\alpha\ne0$ and any finite list containing $\alpha$ whose other
curvatures are homogeneous multiples $b_i\alpha$, with $b_i\in A$.
Then $I=(\alpha)$ and $\mu_p=1$, $\mu_q=0$ for $q\ne p$.
Exactly one curvature handle remains; every other pair splits as in (3).
If $\alpha=0$, all pairs are closed. This holds for every original $X$.
\end{theorem}
\begin{proof}
The image of $\alpha$ generates $I/\m I$. It is nonzero: otherwise
$\alpha=b\alpha$ for some $b\in\m$. Since $b^{n+1}=0$,
$1-b$ has the explicit inverse $1+b+\cdots+b^n$, forcing $\alpha=0$.
The normal-form theorem proves both the sharp lower bound and its
attainment. Thus arbitrarily many mixed-degree births do not dilute the
prescribed input into zero curvature.
\end{proof}

This executes an alpha-dependent construction: form the actual torsor for
$\alpha$, add all chosen cup births, carry out every elimination, and
compute the final canonical base transgression, which is still $\alpha$.
It closes this precise erasure recipe. It does not refute Hodge algebraicity
or rule out another geometric constructor. For $p=1$ an integral multiple
of $\alpha$ is the Chern class of a line bundle and returns its actual
Chow divisor divided by that multiple. In higher degree (1)--(4) alone
do not provide a canonical Chow lift. In particular the ideal algebra
does not assume its Hodge coefficients to be algebraic cycles.

\section{Hostile checks: syzygies are not extra removable handles}
In $A=\Q[x,y]/(x^2,y^2)$ with curvatures $(x,y)$, one has $\mu_1=2$.
The element $y u_1-x u_2$ is closed, since $yx-xy=0$, but all its
coefficients lie in $\m$. It cannot be a row of an invertible homogeneous
matrix: that row vanishes modulo $\m$. Thus a closed syzygy is not a
split free handle. Similarly on $\mathbb P^2$, $h^2u$ is closed for
$du=h$, but $\mu_1=1$ and the curved handle cannot disappear. Nilpotent
annihilators must not be divided out.

For $(h,h^2,h^3)$ on $\mathbb P^3$, the theorem gives one active pair,
two closed pairs of degrees three and five, and
\[
 P_B(t)=(1+t)(1+t^7)(1+t^3)^2(1+t^5)^2.
\]
For $(x,y,xy)$ on $\mathbb P^1\times\mathbb P^1$, it gives two active
pairs and one closed degree-three pair. These are checks of the universal
graded proof, not a replacement special-family objective.

\section{Proof status and formalization boundary}
The finite normal-form, optimality, and cohomology arguments are hand
proofs, with author adversarial self-review. Lean additions encode the
explicit finite simultaneous row clearing, invariant span, nilpotent
fixed-point test, and the nonunit syzygy guard. They do not formalize the
geometric comparison or the complete graded ideal theorem. Compilation
is left to the owner; no kernel receipt or universal Chow proof is claimed.
\end{document}
